\begin{table}[H]\centering
\caption{\textbf{Kitagawa decomposition of the for-profit--public cesarean gap}}
\label{tab:decomposition}
\small
\small\setlength{\tabcolsep}{4pt}\resizebox{\ifdim\width>\linewidth \linewidth\else\width\fi}{!}{%
\begin{tabular}{lcccc}
\toprule
 & \multicolumn{2}{c}{Birth share (\%)} & \multicolumn{2}{c}{Cesarean rate (\%)} \\
\cmidrule(lr){2-3}\cmidrule(lr){4-5}
Robson group & For-profit & Public & For-profit & Public \\
\midrule
1 & 14.6 & 19.8 & 65.8 & 35.1 \\
2 & 22.1 & 10.4 & 89.0 & 55.5 \\
3 & 9.6 & 25.5 & 39.9 & 15.0 \\
4 & 9.4 & 8.8 & 75.6 & 38.0 \\
5 & 28.6 & 19.8 & 94.6 & 78.4 \\
6 & 1.9 & 1.2 & 94.6 & 88.2 \\
7 & 2.3 & 1.9 & 93.6 & 84.8 \\
8 & 2.7 & 2.1 & 93.6 & 80.0 \\
9 & 0.2 & 0.2 & 98.0 & 96.7 \\
10 & 8.5 & 10.3 & 76.9 & 42.8 \\
\midrule
\multicolumn{5}{l}{Total gap: 36.2 percentage points} \\
\multicolumn{5}{l}{\quad from Robson-group composition: 10.0 points (27\%)} \\
\multicolumn{5}{l}{\quad from within-group cesarean rates: 26.3 points (73\%)} \\
\bottomrule
\end{tabular}
}
\begin{minipage}{\linewidth}\footnotesize
\textit{Notes:} SINASC 2014--2024, Robson groups 1--10. Symmetric Kitagawa decomposition of sector-level cesarean rates. Totals are subject to rounding. Within-group differences are descriptive; they do not isolate physician choice or medical inappropriateness.
\end{minipage}
\end{table}
